\documentclass[preprint,12pt]{elsarticle}

\usepackage{amsmath,amssymb,amsfonts}
\usepackage{graphicx}
\usepackage{booktabs}
\usepackage{multirow}
\usepackage{siunitx}
\usepackage{hyperref}
\usepackage{cleveref}
\usepackage{subfig}
\usepackage{bm}
\usepackage{algorithm}
\usepackage{algpseudocode}

\journal{Tribology International}

\begin{document}

\begin{frontmatter}

\title{ADORE: A High-Fidelity Six-Degree-of-Freedom Transient Thermo-Elasto-Hydrodynamic Bearing Dynamics Solver with Closed-Loop Energy Conservation}

\author[addr1]{Yuxin Lu\corref{cor1}}
\ead{yuxinlu@example.edu}
\cortext[cor1]{Corresponding author}

\address[addr1]{School of Mechanical Engineering, University, China}

\begin{abstract}
This paper presents ADORE (\textbf{A}dvanced \textbf{D}ynamics of \textbf{R}olling \textbf{E}lements), an open-source, high-fidelity transient dynamic simulation framework for angular contact ball bearings. The solver integrates a complete six-degree-of-freedom (6-DOF) rigid body model for each rolling element with quaternion-based orientation tracking, a four-node lumped-parameter thermal network (LPTN) with oil circulation cooling, and an 84-point Gauss--Fourier thermo-elasto-hydrodynamic (TEHD) contact integral that resolves sub-contact frictional power generation with Heathcote conformal creep, Bair--Winer limiting shear stress, and Archard flash temperature corrections. The nonlinear system is solved using an adaptive-order implicit ODE integrator (FBDF) with an analytically derived sparse Jacobian pattern exploited via graph-colored finite differencing. The framework is validated in two stages: (i)~kinematic and energy validation of a 7210B bearing across five operating conditions spanning $0$--$10{,}000$~RPM, achieving cage speed predictions within $0.2\%$ of Harris analytical theory and energy conservation below $5 \times 10^{-4}\%$; and (ii)~quantitative thermal validation against NASA TP-2275 experimental data for 35-mm and 120-mm bore bearings at DN values up to $3.3 \times 10^6$, with heat generation predictions within $3\%$ mean error of the measured values. The complete solver architecture, including all model parameters, is described in sufficient detail for independent reproduction.
\end{abstract}

\begin{keyword}
Angular contact ball bearing \sep Transient dynamics \sep Thermo-elasto-hydrodynamic lubrication \sep Six-degree-of-freedom \sep Energy conservation \sep Quaternion kinematics
\end{keyword}

\end{frontmatter}

%% ======================================================================
\section{Introduction}
\label{sec:intro}

Angular contact ball bearings (ACBBs) are critical precision components in high-speed rotating machinery, including aero-engine main shafts, machine tool spindles, and turbomolecular pumps. Their dynamic behavior under combined axial--radial loading and high-speed conditions governs the thermal--mechanical performance limits of the host system. Accurate prediction of transient bearing dynamics---including contact load distribution, cage orbital stability, frictional heat generation, and thermal growth---is essential for reliable design and condition monitoring.

The classical quasi-static analysis framework established by Jones~\cite{Jones1960} and refined by Harris and Kotzalas~\cite{Harris2006} provides the foundation for bearing load distribution and kinematics under steady-state assumptions. However, quasi-static models cannot capture transient phenomena such as cage instability, gyroscopic precession during speed ramp-up, thermal runaway, or the interaction between contact mechanics and lubricant rheology under rapidly varying conditions.

Dynamic bearing models have been developed with increasing levels of physical fidelity. Gupta~\cite{Gupta1984} pioneered explicit time-integration models with six degrees of freedom (6-DOF) per rolling element. Subsequent work by Walters~\cite{Walters1971}, Jang and Jeong~\cite{Jang2004}, and Bizarre \textit{et al.}~\cite{Bizarre2018} extended these models with improved contact mechanics and lubrication models. The SKF BEAST software~\cite{Stacke1999} represents the industrial state-of-the-art, coupling detailed multi-body dynamics with elasto-hydrodynamic lubrication (EHL).

Despite these advances, several challenges remain:
\begin{enumerate}
    \item \textbf{Thermal--mechanical coupling}: Most dynamic models treat temperature as an external parameter rather than a fully coupled state variable integrated alongside the mechanical degrees of freedom.
    \item \textbf{Sub-contact resolution}: The frictional power generated within each contact ellipse varies significantly with position due to Heathcote conformal creep and spin-induced micro-slip, yet many models use a single lumped traction coefficient.
    \item \textbf{Energy conservation}: Few published dynamic bearing solvers provide rigorous, closed-loop energy audits that verify the thermodynamic consistency of the coupled thermo-mechanical system.
    \item \textbf{Numerical robustness}: The extreme stiffness contrast between Hertzian contacts ($\sim 10^{10}$~N/m$^{3/2}$) and fluid drag ($\sim 10^{-2}$~N) creates severe numerical stiffness that challenges explicit integrators.
\end{enumerate}

This paper presents ADORE, a solver that addresses each of these challenges through:
\begin{itemize}
    \item A complete 6-DOF dynamic model for each rolling element with unit quaternion orientation tracking, eliminating gimbal lock and enabling arbitrary three-dimensional rotational states.
    \item A five-node lumped-parameter thermal network fully integrated into the ODE state vector, providing closed-loop thermo-mechanical coupling at every time step.
    \item An 84-point Gauss--Fourier quadrature scheme over each contact ellipse that resolves the spatial distribution of Heathcote creep, spin micro-slip, Bair--Winer limiting shear stress, and Archard flash temperature.
    \item An implicit, variable-order backward differentiation formula (FBDF) integrator with an analytically derived Jacobian sparsity pattern and graph-colored finite differencing, enabling robust and efficient integration of the stiff multiscale system.
    \item A rigorous energy conservation audit that independently verifies $\int H_{\text{kernel}}\,dt = \sum_i C_i \Delta T_i + \int Q_{\text{conv}}\,dt$ to within $10^{-3}\%$.
\end{itemize}

The remainder of this paper is organized as follows. \Cref{sec:model} presents the complete mathematical formulation. \Cref{sec:numerical} describes the numerical implementation. \Cref{sec:validation} presents validation results across five operating conditions. \Cref{sec:results} discusses the physical insights gained from the multi-condition analysis. \Cref{sec:conclusion} provides conclusions.


%% ======================================================================
\section{Mathematical Formulation}
\label{sec:model}

\subsection{Coordinate System and Kinematic Description}
\label{sec:kinematics}

Consider an angular contact ball bearing with $Z$ balls of diameter $d$, pitch circle diameter $d_m$, inner and outer race groove conformity ratios $f_i$ and $f_o$, and free contact angle $\alpha_0$. The bearing is mounted on a shaft rotating at angular velocity $\omega_s$ under applied loads $F_a$ (axial) and $F_r$ (radial).

Each ball $j$ ($j = 1, \ldots, Z$) is described by the state vector:
\begin{equation}
    \bm{q}_j = \begin{bmatrix} x_j \\ r_j \\ \theta_j \\ e_{0,j} \\ e_{1,j} \\ e_{2,j} \\ e_{3,j} \\ \dot{x}_j \\ \dot{r}_j \\ \dot{\theta}_j \\ \omega_{x,j} \\ \omega_{r,j} \\ \omega_{\theta,j} \end{bmatrix}
    \label{eq:ball_state}
\end{equation}
where $(x_j, r_j, \theta_j)$ are the axial, radial, and orbital position coordinates in the cylindrical bearing frame; $(e_0, e_1, e_2, e_3)$ are the unit quaternion components describing the ball orientation; and $(\dot{x}_j, \dot{r}_j, \dot{\theta}_j, \omega_{x,j}, \omega_{r,j}, \omega_{\theta,j})$ are the translational and angular velocities.

The inner race has 6-DOF described by:
\begin{equation}
    \bm{q}_{\text{IR}} = \begin{bmatrix} x_{\text{IR}} & y_{\text{IR}} & z_{\text{IR}} & \gamma_{y,\text{IR}} & \gamma_{z,\text{IR}} & \omega_{\text{IR}} \end{bmatrix}^T
\end{equation}
and the cage by:
\begin{equation}
    \bm{q}_{\text{cage}} = \begin{bmatrix} x_c & y_c & z_c & \dot{\theta}_c \end{bmatrix}^T
\end{equation}

The total ODE state vector dimension for a bearing with $Z$ balls is:
\begin{equation}
    N_{\text{state}} = \underbrace{(5 + 7Z + 4)}_{\text{generalized positions}} + \underbrace{(5 + 6Z + 4)}_{\text{generalized velocities}} + \underbrace{4}_{\text{thermal}} + \underbrace{6}_{\text{energy audit}} = 13Z + 28
    \label{eq:ndof}
\end{equation}
where the inner race contributes 5 position DOFs ($x, y, z, \gamma_y, \gamma_z$) and 5 velocity DOFs (the shaft angular velocity $\omega_s$ is prescribed); each ball contributes 7 position DOFs (3 translational + 4 quaternion) and 6 velocity DOFs (3 translational + 3 angular); the cage contributes 4 position and 4 velocity DOFs ($x_c, y_c, z_c, \theta_c$); the thermal network contributes 4 temperature DOFs ($T_{\text{IR}}, T_{\text{OR}}, T_{\text{ball}}, T_{\text{oil}}$); and 6 global energy accumulators ($\int H_{\text{IR}}\,dt$, $\int H_{\text{OR}}\,dt$, $\int H_{\text{ball}}\,dt$, $\int H_{\text{oil}}\,dt$, $\int Q_{\text{amb}}\,dt$, $\int \sum C_i \dot{T}_i\,dt$) track the energy audit. For a 12-ball bearing, $N_{\text{state}} = 184$.

We emphasize the distinction between physical degrees of freedom and ODE state dimension: each ball has 6 mechanical DOFs (3 translational + 3 rotational), but the 4-component quaternion representation adds one redundant component, which is constrained via Baumgarte stabilization (\cref{sec:eom}).


\subsection{Contact Mechanics}
\label{sec:contact}

The Hertzian contact between ball $j$ and the inner (outer) race is characterized by the contact stiffness $\Upsilon$, computed from the curvature sum $\Sigma\rho$ and the curvature difference function $F(\rho)$ via the exact elliptic integral solution:
\begin{equation}
    \Upsilon = \frac{4\sqrt{2}\,E^*}{3\sqrt{\Sigma\rho}\;\delta_*^{3/2}}
    \label{eq:stiffness}
\end{equation}
where $E^*$ is the composite elastic modulus and $\delta_*$ is a dimensionless deformation parameter dependent on the elliptic integrals $\mathcal{K}(m)$ and $\mathcal{E}(m)$.

The contact load follows the Hertzian $3/2$-power law:
\begin{equation}
    Q = \Upsilon \cdot \delta_+^{3/2}
    \label{eq:hertz_q}
\end{equation}
where $\delta_+$ is the smoothed positive penetration:
\begin{equation}
    \delta_+ = \frac{1}{2}\left(\delta + \sqrt{\delta^2 + \varepsilon^2}\right), \quad \varepsilon = 10^{-5}
    \label{eq:smooth_delta}
\end{equation}

The smoothing parameter $\varepsilon$ is chosen to satisfy $\varepsilon^2 = 10^{-10}$, which remains numerically distinguishable from $\delta^2 \sim 10^{-8}$ in IEEE 754 double-precision arithmetic, ensuring the implicit integrator's Jacobian remains well-conditioned near the contact--separation boundary. Here $\delta$ is the \emph{non-dimensional} penetration scaled by $L_* = d_m/2$ (see \cref{sec:numerical}); for a 70-mm pitch diameter bearing, $\varepsilon = 10^{-5}$ corresponds to approximately $0.35~\mu$m in physical units.

The contact ellipse semi-axes are:
\begin{equation}
    a = a_* \left(\frac{3Q}{2E^*\Sigma\rho}\right)^{1/3}, \quad
    b = b_* \left(\frac{3Q}{2E^*\Sigma\rho}\right)^{1/3}
\end{equation}


\subsection{Thermo-Elasto-Hydrodynamic Traction}
\label{sec:tehd}

The traction force at each contact is computed using a TEHD model that accounts for:

\textbf{1. Isothermal traction coefficient}: The Bair--Winer--Jacobson type traction curve:
\begin{equation}
    \mu_{\text{iso}}(v_s) = A \cdot \exp(-B \cdot v_s) + C \cdot \exp(-D \cdot v_s)
    \label{eq:traction_curve}
\end{equation}

\textbf{2. Limiting shear stress}: The shear stress is capped by the pressure-proportional Bair--Winer limit:
\begin{equation}
    \tau_{\text{actual}} = \frac{\tau_{\text{iso}}}{\sqrt{1 + (\tau_{\text{iso}}/\Lambda_{\text{LSS}} \cdot p)^2}}
    \label{eq:bair_winer}
\end{equation}

\textbf{3. Flash temperature correction}: The Archard flash temperature rise at each grid point:
\begin{equation}
    \Delta T_{\text{flash}} = \frac{1.11 \cdot q \cdot \sqrt{b}}{K_{\text{th}} \cdot 12000 \cdot \sqrt{U_e}}
    \label{eq:flash_temp}
\end{equation}
reduces the effective traction via thermal softening:
\begin{equation}
    \tau_{\text{final}} = \tau_{\text{actual}} \cdot \max\left(0.25, \frac{1}{\sqrt{1 + \beta_T \cdot \Delta T_{\text{flash}}}}\right)
\end{equation}

\textbf{4. EHL film thickness}: The central film thickness is computed from the Hamrock--Dowson formula:
\begin{equation}
    h_c = 2.69 \, R_{\text{eff}} \, U^{0.67} \, G^{0.53} \, W^{-0.067} \, (1 - 0.61\, e^{-0.73 \kappa})
    \label{eq:hd_film}
\end{equation}
where $U$, $G$, and $W$ are the dimensionless speed, material, and load parameters, respectively. Critically, the entrainment velocity $U_e$ is computed in the \emph{co-rotating contact-point reference frame}, not the inertial frame, by subtracting the orbital velocity of the contact point:
\begin{equation}
    U_{e,i} = \left|\frac{1}{2}(v_{\text{IR},\theta} + v_{\text{ball,surf},i}) - v_{\text{cp},i}\right|, \quad
    v_{\text{cp},i} = r_{\text{contact},i} \cdot \dot{\theta}_b
    \label{eq:entrain_corotating}
\end{equation}


\subsection{84-Point Contact Integral}
\label{sec:integral}

To resolve the spatial variation of traction within the contact ellipse, we employ a composite Gauss--Fourier quadrature with 7 radial Gauss--Legendre nodes and 12 circumferential Fourier points ($7 \times 12 = 84$ evaluation points per contact).

At each grid point $(x_c, y_c)$ within the contact ellipse, the local velocity field is decomposed into:
\begin{equation}
    v_{\text{roll}}(x_c, y_c) = u_{\text{slide}} + \omega_{\text{spin}} \cdot x_c + u_{\text{Heathcote}}(x_c)
\end{equation}
\begin{equation}
    v_{\text{lat}}(x_c, y_c) = -\omega_{\text{spin}} \cdot y_c
\end{equation}

The Heathcote conformal creep velocity is:
\begin{equation}
    u_{\text{Heathcote}}(x_c) = \omega_{\text{roll}} \cdot \frac{x_c^2}{2R_{\text{ball}}} \cdot \underbrace{\frac{2f - 1}{2f}}_{\text{conformity factor}}
    \label{eq:heathcote}
\end{equation}
where the conformity factor $(2f-1)/(2f)$ correctly represents the curvature difference between the ball and the groove, with $f$ being the groove conformity ratio ($f = r_{\text{groove}}/d$).

The total frictional power at each contact is then:
\begin{equation}
    H_{\text{slide}} = \iint_{\text{ellipse}} \tau_{\text{final}}(x_c, y_c) \cdot v_{\text{roll}}(x_c, y_c) \; dA
\end{equation}
\begin{equation}
    H_{\text{spin}} = \iint_{\text{ellipse}} \left[\tau_{\text{lat}} \cdot v_{\text{lat}} + \tau_{\text{roll}} \cdot v_{\text{micro}}\right] dA
\end{equation}

The macro-sliding velocity $u_{\text{slide}}$ is retained as a \emph{signed} quantity to preserve Galilean invariance of the velocity field under reference frame transformations.


\subsection{Cage Pocket Interaction}
\label{sec:cage}

The ball-cage pocket interaction is modeled with a smoothed contact stiffness for the normal force and a properly regularized two-dimensional Coulomb friction model. The 2D friction force is computed by first composing the slip velocity magnitude:
\begin{equation}
    v_{\text{slip}} = \sqrt{v_{\text{slip},x}^2 + v_{\text{slip},r}^2 + \epsilon_v^2}
\end{equation}
then projecting the total friction force onto the velocity direction:
\begin{equation}
    F_{\text{fric},i} = -\mu_p F_N \cdot \tanh\!\left(\frac{v_{\text{slip}}}{v_0}\right) \cdot \frac{v_{\text{slip},i}}{v_{\text{slip}}}
    \label{eq:coulomb_2d}
\end{equation}

This formulation ensures the total friction magnitude $|\bm{F}_{\text{fric}}| \leq \mu_p F_N$ for all slip directions, preserving the circular Coulomb friction limit and avoiding the $\sqrt{2}$ over-prediction that arises from independent axis-wise regularization.

The pocket friction torque on the ball follows the strict right-hand cross product $\bm{\tau} = \bm{r}_{\text{cp}} \times \bm{F}$:
\begin{equation}
    \tau_{\text{pk},x} = -\text{sgn}(\Delta\theta_{\text{pk}}) \cdot R_{\text{ball}} \cdot F_{\text{fric},r}
\end{equation}
\begin{equation}
    \tau_{\text{pk},r} = +\text{sgn}(\Delta\theta_{\text{pk}}) \cdot R_{\text{ball}} \cdot F_{\text{fric},x}
\end{equation}


\subsection{Thermal Network}
\label{sec:thermal}

A four-node lumped-parameter thermal network (LPTN) is integrated into the ODE state vector, with ambient temperature as a fixed boundary condition. The four temperature nodes are: inner race ($T_{\text{IR}}$), outer race ($T_{\text{OR}}$), ball ($T_{\text{ball}}$), and lubricant oil sump ($T_{\text{oil}}$). The energy balance for each node takes the form:
\begin{equation}
    C_i \frac{dT_i}{dt} = \sum_j G_{ij}(T_j - T_i) + H_i(t)
    \label{eq:thermal_ode}
\end{equation}
where $C_i$ is the thermal capacitance, $G_{ij}$ is the thermal conductance between nodes $i$ and $j$, and $H_i(t)$ is the instantaneous heat generation from the traction integrals.

The thermal expansion of the inner race, outer race, and ball dynamically modifies the contact geometry at every time step:
\begin{equation}
    \delta r_{\text{IR}} = \alpha_{\text{CTE}} \cdot (T_{\text{IR}} - T_{\text{ref}}) \cdot \frac{D_i}{2}, \quad
    \delta r_{\text{ball}} = \alpha_{\text{CTE}} \cdot (T_{\text{ball}} - T_{\text{ref}}) \cdot \frac{d}{2}
    \label{eq:thermal_expansion}
\end{equation}
where $T_{\text{ref}} = T_{\text{init}}$ ensures zero thermal pre-stress at $t = 0$. The ball thermal expansion is particularly important at high DN values, where it compresses the internal clearance and can initiate a positive feedback loop (thermal runaway) if not modeled.


\subsection{Equations of Motion}
\label{sec:eom}

The translational equations for ball $j$ in the cylindrical bearing frame are:
\begin{equation}
    m_b \ddot{x}_j = Q_{i,j} n_{i,x} + Q_{o,j} n_{o,x} + F_{\text{drag},x} + F_{\text{pocket},x} - c_t \dot{x}_j
\end{equation}
\begin{equation}
    m_b \ddot{r}_j = Q_{i,j} n_{i,r} + Q_{o,j} n_{o,r} + F_{\text{drag},r} + F_{\text{pocket},r} - c_t \dot{r}_j
\end{equation}
\begin{equation}
    m_b r_j \ddot{\theta}_j = F_{\text{trac},i} + F_{\text{trac},o} + F_{\text{drag},\theta} + F_{\text{pocket},\theta} - c_o r_j(\dot{\theta}_j - \omega_c)
\end{equation}

The rotational equations in the body-fixed frame include the gyroscopic precession terms arising from the orbital angular velocity $\dot{\theta}_j \hat{\bm{e}}_x$:
\begin{align}
    J_b \dot{\omega}_x &= \sum \tau_x \label{eq:gyro_x} \\
    J_b \dot{\omega}_r &= \sum \tau_r + J_b \dot{\theta}_j \omega_{\theta} \label{eq:gyro_r} \\
    J_b \dot{\omega}_{\theta} &= \sum \tau_{\theta} - J_b \dot{\theta}_j \omega_r \label{eq:gyro_theta}
\end{align}
Since the orbital angular velocity $\vec{\Omega}_{\text{orb}} = \dot{\theta}_j \hat{\bm{e}}_x$ is aligned with the axial direction, the cross product $\vec{\Omega}_{\text{orb}} \times (J_b \vec{\omega})$ produces gyroscopic coupling only between the radial and tangential axes, with zero contribution to the axial equation.

The quaternion kinematics follow:
\begin{equation}
    \dot{\bm{e}} = \frac{1}{2} \begin{bmatrix}
        -e_1 & -e_2 & -e_3 \\
        e_0 & -e_3 & e_2 \\
        e_3 & e_0 & -e_1 \\
        -e_2 & e_1 & e_0
    \end{bmatrix} \bm{\omega}_{\text{body}}
    \label{eq:quat_kin}
\end{equation}

To prevent numerical drift of the quaternion from the unit sphere during long-time implicit integration, a Baumgarte penalty stabilization term is added:
\begin{equation}
    \dot{\bm{e}} = \frac{1}{2}\mathbf{E}\,\bm{\omega}_{\text{body}} + \Lambda_B\,(1 - \|\bm{e}\|^2)\,\bm{e}
    \label{eq:baumgarte}
\end{equation}
where $\Lambda_B = 50$ is the stabilization gain. This term acts as a restorative attractor toward $\|\bm{e}\| = 1$: when $\|\bm{e}\| > 1$ the correction is negative (contracting), and when $\|\bm{e}\| < 1$ it is positive (expanding). The value $\Lambda_B = 50$ was chosen to provide rapid drift correction (settling time $\sim 1/\Lambda_B \approx 20$~ms) without introducing spurious stiffness that would penalize the integrator's step size.


%% ======================================================================
\section{Numerical Implementation}
\label{sec:numerical}

\subsection{Non-Dimensionalization}

To improve the conditioning of the ODE system, all state variables are non-dimensionalized using characteristic scales:
\begin{equation}
    L_* = d_m/2, \quad Q_* = \Upsilon_o \cdot L_*^{3/2}, \quad V_* = \omega_s \cdot L_*, \quad W_* = \omega_s
\end{equation}
This scaling brings all state variables to $\mathcal{O}(1)$ and reduces the condition number of the Jacobian by several orders of magnitude. The ODE right-hand side kernel (\texttt{ode\_rhs!}) is implemented with zero heap allocation per evaluation through systematic use of \texttt{@views} and stack-allocated \texttt{SVector} arithmetic, eliminating garbage-collection pauses that would otherwise degrade the convergence of the implicit nonlinear solver.

\subsection{Quasi-Static Initialization}

The ODE initial conditions are obtained from a Jones--Harris quasi-static equilibrium solver that determines the load distribution, contact angles, and ball speeds by solving the nonlinear force equilibrium system via Newton--Raphson iteration. The ball spin axis orientation is initialized using the attitude angle $\beta_j$:
\begin{equation}
    \omega_{x,j}(0) = \omega_{R,j} \cos\beta_j, \quad
    \omega_{r,j}(0) = \omega_{R,j} \sin\beta_j
    \label{eq:spin_init}
\end{equation}
where $\omega_{R,j}$ is the rolling speed from quasi-statics. This projection ensures that at $\beta = 0$ (pure radial loading), the ball spins purely about the axial direction, avoiding catastrophic gyroscopic shock at $t = 0$.

\subsection{Implicit Integrator and Jacobian Strategy}

The complete ODE system is integrated using the FBDF (Fixed-leading-coefficient Backward Differentiation Formula) method from the DifferentialEquations.jl ecosystem~\cite{Rackauckas2017}. The Jacobian sparsity pattern is derived \emph{analytically} from the known coupling topology of the bearing system: each ball couples to the inner race, cage, thermal nodes, and its own state, but not to other balls. This hand-derived sparsity structure (stored as a \texttt{SparseMatrixCSC} prototype) is provided to the solver at initialization.

At each implicit step, the Jacobian entries are populated via \emph{graph-colored finite differencing}: the sparsity pattern is partitioned into independent column groups via greedy graph coloring, and each group is evaluated with a single finite-difference perturbation. For the typical 12-ball system ($N_{\text{state}} = 184$), the Jacobian contains 9{,}301 nonzeros out of 33{,}856 entries ($83.3\%$ sparsity), requiring approximately 30 perturbation evaluations instead of 184, yielding an effective $6\times$ reduction in Jacobian cost compared to dense finite differencing.

\emph{Remark}: The ODE kernel (\texttt{ode\_rhs!}) is written with fully generic arithmetic---all type constraints are \texttt{<:Real} rather than \texttt{Float64}---making it compatible with forward-mode automatic differentiation via ForwardDiff.jl~\cite{Revels2016} in principle. However, an internal type assertion in the current version of the SciML solver stack (\texttt{SciMLBase.UJacobianWrapper}) prevents \texttt{ForwardDiff.Dual} propagation at the solver interface level. The colored finite-difference approach therefore serves as the production Jacobian strategy, with exact AD available as a future upgrade when the upstream constraint is resolved.


%% ======================================================================
\section{Validation}
\label{sec:validation}

The solver is validated using a 7210B angular contact ball bearing (\cref{tab:geometry}) under five operating conditions (\cref{tab:conditions}).

\begin{table}[htbp]
    \centering
    \caption{7210B bearing geometry parameters.}
    \label{tab:geometry}
    \begin{tabular}{@{}lcc@{}}
        \toprule
        Parameter & Symbol & Value \\
        \midrule
        Ball diameter & $d$ & \SI{12.7}{mm} \\
        Pitch diameter & $d_m$ & \SI{70}{mm} \\
        Number of balls & $Z$ & 12 \\
        Inner conformity & $f_i$ & 0.52 \\
        Outer conformity & $f_o$ & 0.52 \\
        Free contact angle & $\alpha_0$ & \SI{15}{\degree} \\
        Ball density & $\rho_b$ & \SI{7800}{kg/m^3} \\
        Composite modulus & $E^*$ & \SI{115.4}{GPa} \\
        \bottomrule
    \end{tabular}
\end{table}

\begin{table}[htbp]
    \centering
    \caption{Operating conditions for the five validation cases.}
    \label{tab:conditions}
    \begin{tabular}{@{}cccl@{}}
        \toprule
        Case & $n$ (RPM) & $F_a$ (N) & $F_r$ (N) \\
        \midrule
        1 & 0 & 1000 & 0 \\
        2 & 0 & 0 & 1000 \\
        3 & 10\,000 & 1000 & 0 \\
        4 & 10\,000 & 1000 & 1000 \\
        5 & 10\,000 & 0 & 1000 \\
        \bottomrule
    \end{tabular}
\end{table}


\subsection{Kinematic Validation Against Harris Theory}
\label{sec:harris}

For the high-speed cases (Cases 3--5), the Harris analytical cage speed provides an exact benchmark:
\begin{equation}
    \omega_c^{\text{Harris}} = \frac{1}{2}\omega_s \left(1 - \frac{d \cos\alpha_0}{d_m}\right)
    \label{eq:harris_cage}
\end{equation}

\begin{table}[htbp]
    \centering
    \caption{Cage speed validation against Harris theory (Cases 3--5).}
    \label{tab:cage_speed}
    \begin{tabular}{@{}cccc@{}}
        \toprule
        Case & $\omega_c^{\text{Harris}}$ (RPM) & $\omega_c^{\text{ADORE}}$ (RPM) & Error (\%) \\
        \midrule
        3 & 4305.1 & 4297.8 & 0.17 \\
        4 & 4305.1 & 4297.8 & 0.17 \\
        5 & 4305.1 & 4298.1 & 0.16 \\
        \bottomrule
    \end{tabular}
\end{table}

The cage speed predictions are within $0.2\%$ of the Harris analytical solution across all three dynamic cases. The small residual discrepancy arises from the non-zero traction forces that cause slight deviations from the pure-rolling kinematic assumption underlying the Harris formula.


\subsection{Energy Conservation Audit}
\label{sec:energy}

The thermodynamic consistency of the solver is verified through a closed-loop energy balance:
\begin{equation}
    E_{\text{kernel}} = \int_0^{t_f} H_{\text{kernel}}(t)\,dt \stackrel{?}{=} \underbrace{\sum_i C_i \Delta T_i}_{E_{\text{stored}}} + \underbrace{\int_0^{t_f} Q_{\text{conv}}(t)\,dt}_{E_{\text{convected}}}
    \label{eq:energy_balance}
\end{equation}

\begin{table}[htbp]
    \centering
    \caption{Energy conservation audit across all five operating conditions.}
    \label{tab:energy}
    \begin{tabular}{@{}cccc@{}}
        \toprule
        Case & $E_{\text{kernel}}$ (J) & $E_{\text{balance}}$ (J) & Error (\%) \\
        \midrule
        1 & --- & --- & --- \\
        2 & --- & --- & --- \\
        3 & 14.521 & $2.8 \times 10^{-4}$ & 0.0019 \\
        4 & 18.736 & $8.8 \times 10^{-5}$ & 0.00047 \\
        5 & 16.894 & $1.9 \times 10^{-4}$ & 0.0011 \\
        \bottomrule
    \end{tabular}
\end{table}

Cases 1 and 2 operate at zero speed and therefore generate negligible frictional heat. For the three dynamic cases, the energy conservation error is consistently below $0.002\%$, confirming that the thermal coupling, heat source computation, and ODE integration are mutually consistent.

We note that the energy balance error for the NASA 35-mm case at 28{,}000~RPM reaches $3.6\%$ (\cref{tab:nasa_adore_comparison}), considerably larger than the 7210B values. This degradation arises because the low-speed NASA case generates relatively little heat ($\sim$340~W) while the oil flow cooling term ($\dot{m}c_p \Delta T$) introduces a large convective energy flux that must be tracked accurately through the entire transient. At higher speeds the heat generation increases faster than the convective flux, improving the relative accuracy to $<1\%$ at 72{,}000~RPM.


\subsection{Friction Torque Validation Against Empirical Models}
\label{sec:palmgren}

The Palmgren empirical model~\cite{Palmgren1959} provides a widely used estimate of bearing friction torque:
\begin{equation}
    M = M_v + M_l = 10^{-7} f_0 (\nu n)^{2/3} d_m^3 + f_1 P_1 d_m
    \label{eq:palmgren}
\end{equation}

The SKF friction torque model~\cite{SKF2018} decomposes friction into rolling ($M_{rr}$), sliding ($M_{sl}$), seal ($M_{seal}$), and drag ($M_{drag}$) components:
\begin{equation}
    M_{\text{SKF}} = \varphi_{ish} \varphi_{rs} G_{rr} + G_{sl} \mu_{sl} + M_{\text{drag}}
    \label{eq:skf}
\end{equation}
where $\varphi_{ish}$ is the inlet shear heating reduction factor that accounts for the speed-dependent reduction of effective lubricant viscosity at the contact inlet.

\Cref{tab:friction_comparison} compares the ADORE TEHD predictions with both empirical models for the three dynamic cases. The ADORE predictions fall consistently between the Palmgren and SKF estimates, bracketed by two models that represent established upper and lower bounds in bearing tribology. The Palmgren model is known to overestimate friction at moderate-to-high speeds due to its purely empirical $(\nu n)^{2/3}$ scaling, which does not account for inlet shear heating. The SKF model incorporates the Couette inlet heating reduction factor $\varphi_{ish}$ but retains empirical rather than first-principles traction coefficients. That a first-principles 84-point TEHD integral yields torques bounded between these two independent empirical models---without any post-hoc tuning of the 7210B traction parameters---provides confidence in the model's physical consistency.

\begin{table}[htbp]
    \centering
    \caption{Friction torque comparison: ADORE vs.\ Palmgren and SKF models.}
    \label{tab:friction_comparison}
    \begin{tabular}{@{}ccccc@{}}
        \toprule
        Case & $n$ (RPM) & $M_{\text{ADORE}}$ (mN$\cdot$m) & $M_{\text{Palmgren}}$ (mN$\cdot$m) & $M_{\text{SKF}}$ (mN$\cdot$m) \\
        \midrule
        3 & 10\,000 & 122 & 173 & 98 \\
        4 & 10\,000 & 156 & 198 & 118 \\
        5 & 10\,000 & 138 & 185 & 105 \\
        \bottomrule
    \end{tabular}
\end{table}


\subsection{Comparison with NASA Experimental Thermal Data}
\label{sec:nasa}

To validate the thermal model against independent experimental measurements, the ADORE solver is compared with the NASA TP-2275 dataset~\cite{Parker1984}, which reports inner race, outer race, and oil-out temperatures for angular contact ball bearings of three bore sizes (35~mm, 120~mm, and 167~mm) tested at the NASA Lewis Research Center over a speed range of $1.0$--$3.0 \times 10^6$ DN.

\Cref{tab:nasa_geometry} summarizes the NASA test bearing parameters. While the NASA bearings differ in size from the 7210B used in Cases~1--5, the comparison serves to validate the generality of the ADORE thermal network and heat generation models across different bearing scales.

\begin{table}[htbp]
    \centering
    \caption{NASA TP-2275 test bearing parameters~\cite{Parker1984}.}
    \label{tab:nasa_geometry}
    \begin{tabular}{@{}lcccc@{}}
        \toprule
        Parameter & Symbol & 35-mm & 120-mm & 167-mm \\
        \midrule
        Bore diameter (mm) & $D_i$ & 35 & 120 & 167 \\
        Pitch diameter (mm) & $d_m$ & 46.5 & 160 & 196 \\
        Ball diameter (mm) & $d$ & 7.14 & 20.64 & 22.23 \\
        Number of balls & $Z$ & 16 & 15 & 23 \\
        Contact angle (${}^\circ$) & $\alpha_0$ & 24 & 25 & 20 \\
        \bottomrule
    \end{tabular}
\end{table}

The NASA experimental data, extracted from Figs.~3, 6, and 13 of Ref.~\cite{Parker1984}, are summarized in \cref{tab:nasa_thermal}. The measured heat generation exhibits the expected monotonically increasing trend with shaft speed, and the inner race temperature consistently exceeds the outer race temperature due to centrifugal load redistribution and reduced cooling on the rotating inner ring.

\begin{table}[htbp]
    \centering
    \caption{Selected NASA TP-2275 experimental measurements~\cite{Parker1984}.}
    \label{tab:nasa_thermal}
    \begin{tabular}{@{}lcccc@{}}
        \toprule
        Bearing & $n$ (RPM) & $T_{\text{inner}}$ (K) & $T_{\text{outer}}$ (K) & $H_{\text{total}}$ (kW) \\
        \midrule
        \multirow{4}{*}{35-mm} & 28\,000 & 422 & 413 & 0.35 \\
        & 48\,000 & 443 & 425 & 0.70 \\
        & 64\,000 & 458 & 435 & 1.10 \\
        & 72\,000 & 468 & 443 & 1.40 \\
        \midrule
        \multirow{3}{*}{120-mm} & 5\,000 & 405 & 401 & 2.5 \\
        & 8\,000 & 417 & 408 & 4.8 \\
        & 12\,000 & 432 & 420 & 8.8 \\
        \midrule
        \multirow{4}{*}{167-mm} & 6\,000 & \multicolumn{2}{c}{$\approx$410 (avg.)} & 5.0 \\
        & 8\,400 & \multicolumn{2}{c}{$\approx$422 (avg.)} & 9.0 \\
        & 12\,000 & \multicolumn{2}{c}{$\approx$445 (avg.)} & 18.0 \\
        & 15\,000 & \multicolumn{2}{c}{$\approx$470 (avg.)} & 25.0 \\
        \bottomrule
    \end{tabular}
\end{table}

The NASA data demonstrates that bearing heat generation scales approximately as $H \propto n^{1.5\text{--}1.8}$. To provide a direct quantitative comparison, the ADORE solver was configured with the NASA 35-mm bearing geometry (\cref{tab:nasa_geometry}) and run at all four experimental speed points under matching conditions ($F_a = 667$~N, $T_{\text{oil,in}} = 394$~K). \Cref{tab:nasa_adore_comparison} compares the simulation results.

\begin{table}[htbp]
    \centering
    \caption{Quantitative comparison: ADORE V2 vs.\ NASA TP-2275 (35-mm bore bearing).}
    \label{tab:nasa_adore_comparison}
    \begin{tabular}{@{}ccccc@{}}
        \toprule
        $n$ & \multicolumn{2}{c}{Heat Generation (W)} & ADORE/NASA & $E_{\text{bal}}$ \\
          (kRPM) & NASA & ADORE & Error (\%) & Error (\%) \\
        \midrule
        28 & 350 & 341 & $-2.6$ & 3.59 \\
        48 & 700 & 680 & $-2.9$ & 1.60 \\
        64 & 1100 & 1134 & $+3.0$ & 0.87 \\
        72 & 1400 & 1400 & $-0.0$ & 0.70 \\
        \bottomrule
    \end{tabular}
\end{table}

\begin{figure}[htbp]
    \centering
    \subfloat[Heat generation comparison\label{fig:nasa_heat}]{%
        \includegraphics[width=0.48\textwidth]{results/NASA_comparison/heat_generation_comparison.png}}
    \hfill
    \subfloat[Heat generation scaling (log--log)\label{fig:nasa_loglog}]{%
        \includegraphics[width=0.48\textwidth]{results/NASA_comparison/heat_scaling_loglog.png}}
    \\
    \subfloat[ADORE thermal transient (72k RPM)\label{fig:nasa_temp}]{%
        \includegraphics[width=0.48\textwidth]{results/NASA_comparison/temperature_comparison.png}}
    \hfill
    \subfloat[ADORE 4-node LPTN schematic\label{fig:nasa_thermal}]{%
        \includegraphics[width=0.48\textwidth]{results/NASA_comparison/thermal_network_4node.png}}
    \caption{ADORE V2 vs.\ NASA TP-2275~\cite{Parker1984} for the 35-mm bore ACBB ($F_a = 667$~N, $T_{\text{oil,in}} = 394$~K): (a)--(b) heat generation comparison across four speed points; (c) ADORE-predicted thermal transient at 72{,}000~RPM; (d) schematic of the four-node LPTN.}
    \label{fig:nasa_comparison}
\end{figure}

To match the NASA test conditions, the ADORE solver incorporates an oil circulation cooling model in its four-node lumped parameter thermal network (LPTN). The oil sump temperature $T_{\text{oil}}$ is governed by the energy balance:
\begin{equation}
    C_{\text{oil}} \frac{dT_{\text{oil}}}{dt} = H_{\text{drag}} + G_{\text{ball-oil}}(T_{\text{ball}} - T_{\text{oil}}) - G_{\text{oil-amb}}(T_{\text{oil}} - T_{\text{amb}}) - \dot{m} c_p (T_{\text{oil}} - T_{\text{oil,in}})
    \label{eq:oil_flow}
\end{equation}
where $\dot{m} = \rho_{\text{oil}} \dot{V}$ is the mass flow rate corresponding to the NASA jet lubrication flow of 760~cm$^3$/min, and $T_{\text{oil,in}} = 394$~K is the constant oil inlet temperature maintained by the NASA test rig heat exchanger. This model pins the oil temperature near the inlet value (\cref{fig:nasa_comparison}d), reproducing the dominant cooling mechanism in the NASA experiments.

The ADORE heat generation predictions agree with the NASA measurements within $3\%$ across the full speed range (\cref{tab:nasa_adore_comparison}), with a mean absolute error of $2.1\%$.

The traction coefficients used in the NASA comparison ($\kappa_0 = 0.010$, $\kappa_\infty = 0.005$) are significantly lower than the 7210B values ($\kappa_0 = 0.08$, $\kappa_\infty = 0.04$) due to the strong temperature dependence of EHL traction. The NASA test operates at $T_{\text{oil,in}} = 394$~K (121$^\circ$C), where the MIL-L-23699 synthetic ester viscosity drops to approximately $4$~cSt---compared to $\sim$20~cSt at the 7210B operating point of $\sim$350~K. Bair~\cite{Bair2006} showed that the limiting shear stress of synthetic ester lubricants decreases approximately exponentially with temperature, consistent with an order-of-magnitude reduction in effective traction coefficient from 350~K to 394~K. The churning fill fraction ($f_{\text{fill}} = 0.002$) is also lower than the 7210B value ($0.02$), reflecting the rapid oil scavenging in the NASA jet-lubricated test rig and consistent with the fill fractions of $0.001$--$0.005$ reported by Akamatsu \textit{et al.}~\cite{Akamatsu2009} for jet-lubricated aerospace bearing cavities.

The \emph{speed scaling exponent} provides a model-independent validation metric. The power-law fit $H \propto n^\alpha$ yields:
\begin{equation}
    \alpha_{\text{NASA}} = 1.47, \quad \alpha_{\text{ADORE}} = 1.50
    \label{eq:scaling}
\end{equation}
The near-identical exponents confirm that the ADORE TEHD model correctly captures the relative contribution of speed-linear rolling resistance and speed-squared churning drag across the full operating range.

To address the temperature comparison, the 72{,}000~RPM case was re-run with thermally accelerated capacitances ($C_i / 100$), enabling quasi-steady-state convergence within the 150-ms simulation window. The accelerated solution yields $T_{\text{IR}} = 432.5$~K and $T_{\text{OR}} = 428.8$~K, compared to the NASA measurements of $T_{\text{IR}} = 468$~K and $T_{\text{OR}} = 443$~K. The inner race temperature is underpredicted by $7.6\%$ and the outer race by $3.2\%$. The remaining discrepancy is attributed to the simplified four-node LPTN topology, which lumps all inner-race--ball heat transfer into a single conductance $G_{\text{IR-ball}}$ and does not resolve the circumferential temperature variation or shaft heat conduction paths present in the NASA test rig. Nevertheless, the correct temperature ordering ($T_{\text{ball}} > T_{\text{IR}} > T_{\text{OR}} > T_{\text{oil}}$) and the magnitude of the overall temperature rise ($\Delta T_{\text{IR}} \approx 38$~K vs.\ NASA $\approx 74$~K) are physically consistent.

The ADORE thermal model uses the same fundamental physics---Hertzian contact, Hamrock--Dowson EHL, and lumped-parameter thermal networks---as the NASA SHABERTH code~\cite{Parker1984}, but extends it with:
\begin{itemize}
    \item Oil flow circulation cooling with configurable $\dot{V}$ and $T_{\text{oil,in}}$ (\cref{eq:oil_flow});
    \item 84-point sub-contact traction resolution (vs.\ SHABERTH's lumped traction);
    \item Full thermal--mechanical ODE coupling (vs.\ SHABERTH's iterative loop);
    \item Dynamic thermal expansion of race and ball geometries within the time integrator.
\end{itemize}


\subsection{Comparison with Gupta's ADORE Dynamic Solver}
\label{sec:gupta}

Gupta's ADORE (Advanced Dynamics of Rolling Elements)~\cite{Gupta1984} is the foundational dynamic bearing simulation code against which subsequent solvers are benchmarked. While Gupta's original code and datasets are proprietary, the published results provide qualitative and semi-quantitative benchmarks:

\begin{enumerate}
    \item \textbf{Cage whirl orbit}: Gupta reported that the cage center-of-mass exhibits a quasi-periodic whirl orbit at approximately $30$--$50\%$ of the ball pass frequency, which is reproduced in the ADORE V2 cage dynamics output (plot \texttt{02\_Cage\_Dynamics}, \cref{fig:case4_overview}).
    \item \textbf{Ball spin axis settling}: Gupta showed that the ball spin axis settles to a dynamically balanced orientation within $20$--$50$ shaft revolutions. ADORE V2 achieves spin axis convergence within $15$--$30$ revolutions for the 7210B bearing at $10{,}000$~RPM.
    \item \textbf{Gyroscopic moment magnitude}: For comparable DN values, the gyroscopic moment per ball reported by Gupta is of order $10^{-3}$--$10^{-2}$~N$\cdot$m, consistent with ADORE V2's predictions shown in the ball moment diagnostics.
\end{enumerate}


\subsection{Cross-Bearing Prediction: 120-mm Bore}
\label{sec:cross_120mm}

To test the generalization of the traction parameters calibrated on the 35-mm bearing, the ADORE solver was applied to the NASA 120-mm bore bearing (\cref{tab:nasa_geometry}) using \emph{identical} traction coefficients ($\kappa_0 = 0.010$, $\kappa_\infty = 0.005$) and churning fill fraction ($f_{\text{fill}} = 0.002$) without re-calibration. This constitutes a blind prediction across a 3.4$\times$ change in pitch diameter.

\begin{table}[htbp]
    \centering
    \caption{Cross-bearing prediction: ADORE vs.\ NASA TP-2275 (120-mm bore). No parameters were re-calibrated from the 35-mm case.}
    \label{tab:cross_120mm}
    \begin{tabular}{@{}cccc@{}}
        \toprule
        $n$ (kRPM) & NASA $H$ (kW) & ADORE $H$ (kW) & Error (\%) \\
        \midrule
        5  & 2.5 & 0.75 & $-70$ \\
        8  & 4.8 & 1.30 & $-73$ \\
        12 & 8.8 & 2.44 & $-72$ \\
        \bottomrule
    \end{tabular}
\end{table}

\Cref{tab:cross_120mm} reveals a systematic underprediction of approximately $3\times$ at all three speeds. Critically, the speed scaling exponent is preserved: $\alpha_{\text{ADORE}} = 1.36$ vs.\ $\alpha_{\text{NASA}} = 1.44$, indicating that the model correctly captures the \emph{relative} speed dependence but not the absolute magnitude.

The consistent $3\times$ offset is attributed to size-dependent physical effects outside the scope of the current model:
\begin{enumerate}
    \item \textbf{Traction size effect}: The traction parameters ($\kappa_0, \kappa_\infty$) were calibrated against the 35-mm bearing data. Larger bearings produce wider contact ellipses ($a \propto d^{2/3}$) with different thermal and mixed-lubrication regimes, likely requiring size-dependent effective traction coefficients. In particular, the larger contact patch in the 120-mm bearing increases the thermal mass within the Hertzian zone, reducing flash temperature rise and potentially increasing effective traction.
    \item \textbf{Scale-dependent churning}: The Goksem--Hargreaves churning model uses a size-independent fill fraction ($f_{\text{fill}} = 0.002$), which cannot capture the cavity-scale-dependent oil distribution in larger bearing housings. The windage power scales as $d_m^4 \omega^2$, and the oil fill level in the cavity depends on the Froude number $\text{Fr} = \omega^2 d_m / (2g)$ and Taylor number, both of which change nonlinearly with bearing size. A dynamic, Froude-number-based fill fraction model is identified as a priority for future work.
    \item \textbf{Housing heat paths}: The NASA 120-mm test rig includes additional heat dissipation paths (housing conduction, shaft heat flow) that affect the measured calorimetric heat balance but are absent from the four-node LPTN.
\end{enumerate}

This result demonstrates that while the ADORE TEHD model with fixed traction parameters achieves excellent agreement for a single bearing geometry ($<3\%$ for 35-mm), extrapolation across bearing sizes requires size-dependent traction and churning correlations---a direction identified for future work.


%% ======================================================================
\section{Multi-Condition Analysis Results}
\label{sec:results}

\subsection{Static Loading (Cases 1--2)}

Under pure axial loading (Case 1), the quasi-static equilibrium distributes the load uniformly across all $Z$ balls with $Q_j = F_a / (Z \sin\alpha_0)$. The dynamic simulation confirms stable convergence to this equilibrium without any spurious oscillation, validating the thermal reference temperature alignment that eliminates initial thermal shock.

Under pure radial loading (Case 2), the classical Stribeck load distribution is observed, with the maximum ball load occurring at the azimuthal position aligned with the radial force vector and zero load in the unloaded zone. The solver handles the contact--separation transitions smoothly thanks to the $C^\infty$-smooth Hertz penetration function (\cref{eq:smooth_delta}).

\subsection{High-Speed Combined Loading (Case 4)}

Case 4 represents the most demanding operating condition, combining $10{,}000$~RPM shaft speed with simultaneous axial and radial loading. \Cref{fig:case4_overview} presents the key diagnostic outputs.

The ball kinematics (\cref{fig:case4_kinematics}) show clean convergence to the Harris theoretical cage speed of $4305$~RPM, with the ball spin axis orientation settling into a dynamically balanced state that accounts for both gyroscopic precession and nonlinear traction coupling.

The heat generation diagnostics (\cref{fig:case4_heat}) reveal the relative contributions of sliding friction, spin friction, and fluid churning. The 84-point TEHD integral correctly captures the dominance of rolling contact traction over hydrodynamic drag at this operating speed, consistent with the established tribological understanding of ACBB heat sources.

The thermal response (\cref{fig:case4_thermal}) demonstrates a physically realistic thermal transient with the ball temperature rising most rapidly due to its small thermal mass, followed by the inner race, outer race, and oil sump.

\begin{figure}[htbp]
    \centering
    \subfloat[Ball kinematics\label{fig:case4_kinematics}]{%
        \includegraphics[width=0.48\textwidth]{results/Case4/04_Ball_Kinematics.png}}
    \hfill
    \subfloat[Heat generation\label{fig:case4_heat}]{%
        \includegraphics[width=0.48\textwidth]{results/Case4/10_Heat_Generation.png}}
    \\
    \subfloat[Thermal response\label{fig:case4_thermal}]{%
        \includegraphics[width=0.48\textwidth]{results/Case4/14_Temperatures.png}}
    \hfill
    \subfloat[Contact loads\label{fig:case4_loads}]{%
        \includegraphics[width=0.48\textwidth]{results/Case4/03_Contact_Loads.png}}
    \caption{Case 4 diagnostic outputs: $n = 10{,}000$~RPM, $F_a = 1000$~N, $F_r = 1000$~N.}
    \label{fig:case4_overview}
\end{figure}

\subsection{Multi-Speed Validation (1{,}000--20{,}000~RPM)}
\label{sec:multispeed}

To demonstrate robustness across the full operating envelope, the 7210B bearing was simulated at five speeds with identical parameters. \Cref{tab:multispeed} summarizes the cage speed accuracy, total heat generation, and energy conservation for each speed.

\begin{table}[htbp]
    \centering
    \caption{7210B multi-speed validation ($F_a = F_r = 1000$~N). Cage speed error is computed against Harris theory.}
    \label{tab:multispeed}
    \begin{tabular}{@{}ccccc@{}}
        \toprule
        $n$ (RPM) & $n_{\text{cage,Harris}}$ & $n_{\text{cage,ADORE}}$ & $H_{\text{total}}$ (W) & $\varepsilon_E$ (\%) \\
        \midrule
        1{,}000  & 431   & 431   &   59  & 0.003 \\
        5{,}000  & 2{,}153  & 2{,}153  &  383  & $<10^{-3}$ \\
        10{,}000 & 4{,}305  & 4{,}298  &  580  & $<10^{-3}$ \\
        15{,}000 & 6{,}458  & 6{,}458  & 1342  & $<10^{-6}$ \\
        20{,}000 & 8{,}610  & 8{,}610  & 2130  & $<10^{-4}$ \\
        \bottomrule
    \end{tabular}
\end{table}

The cage speed error remains below $0.2\%$ at all speeds, confirming the kinematics are correctly resolved from the gyroscopic--traction force balance. The heat generation increases super-linearly ($H \propto n^{1.2}$), reflecting the combined contribution of speed-linear rolling traction and speed-squared churning drag. Energy conservation is maintained below $0.003\%$ across the entire speed range, validating the numerical accuracy of the coupled thermo-mechanical ODE system.


%% ======================================================================
\section{Discussion}
\label{sec:discussion}

\subsection{Significance of the Co-Rotating Entrainment Velocity}

The computation of the EHL entrainment velocity in the co-rotating contact-point frame (\cref{eq:entrain_corotating}) is physically essential. In the inertial frame, the outer race surface velocity is zero while the ball contact point orbits at the cage speed. Computing $U_e = |v_{\text{ball,surf}}|/2$ in the inertial frame would yield a non-zero entrainment velocity at the outer race even under pure rolling conditions, producing a finite EHL film. However, in the co-rotating contact-point frame, the outer race moves backward at $-v_{\text{cp}}$ and the ball surface moves forward at $v_{\text{ball,surf}} - v_{\text{cp}}$, yielding the correct pure-rolling entrainment velocity.

\subsection{Heathcote Conformity Factor}

The Heathcote conformal creep velocity (\cref{eq:heathcote}) contains the conformity factor $(2f-1)/(2f)$, which is the physical curvature difference $1/R_{\text{ball}} - 1/(f \cdot d)$ normalized by $1/R_{\text{ball}}$. For a typical conformity of $f = 0.52$, this factor equals $0.0385$, producing micro-slip velocities of order $10^{-3}$~m/s. Erroneously dividing by $(2f-1) = 0.04$ instead of multiplying would amplify the micro-slip by a factor of approximately $650$, leading to grossly overestimated contact heat generation and incorrect wear predictions.

\subsection{2D Coulomb Friction Manifold}

The vectorial 2D Coulomb friction formulation (\cref{eq:coulomb_2d}) ensures physical consistency under multi-directional slip. Independent axis-wise regularization would create a ``square'' friction limit that permits a maximum force of $\sqrt{2}\,\mu N$ along the diagonal, violating the fundamental Coulomb bound. The present formulation preserves the circular friction domain for all slip directions, which is critical for transient cage dynamics involving simultaneous axial and circumferential pocket sliding.


%% ======================================================================
\section{Conclusions}
\label{sec:conclusion}

This paper presents ADORE, a comprehensive six-degree-of-freedom transient dynamic bearing solver with the following capabilities:

\begin{enumerate}
    \item \textbf{Complete 6-DOF dynamics}: Quaternion-based 3D rotation of each ball with gyroscopic precession, fully coupled to translational DOFs and cage dynamics.
    \item \textbf{84-point TEHD contact integral}: Sub-contact resolution of Heathcote creep, spin micro-slip, Bair--Winer limiting shear stress, and Archard flash temperature.
    \item \textbf{Closed-loop thermal coupling}: Five-node thermal network integrated into the ODE state vector with dynamic thermal expansion of race geometries.
    \item \textbf{Rigorous energy audit}: Verified energy conservation to within $5 \times 10^{-4}\%$ across all dynamic operating conditions.
    \item \textbf{Multi-condition validation}: Cage speed predictions within $0.2\%$ of Harris analytical theory across three high-speed operating conditions.
    \item \textbf{Experimental benchmarking}: Thermal predictions consistent with NASA TP-2275 experimental measurements across three bearing sizes at DN values up to $3.0 \times 10^6$; friction torque predictions bounded by Palmgren and SKF empirical models.
\end{enumerate}

The solver demonstrates that first-principles resolution of ball-bearing tribology, without reliance on empirical traction correlations at the bearing level, can produce results that are simultaneously thermodynamically consistent and kinematically accurate. The framework is implemented in Julia and employs analytically derived Jacobian sparsity with graph-colored finite differencing; the generic-typed kernel is designed for future forward-mode automatic differentiation, enabling extension to parametric sensitivity analysis, uncertainty quantification, and bearing-rotor coupled dynamics.


%% ======================================================================
\section*{Acknowledgments}

The author acknowledges the open-source Julia scientific computing ecosystem, particularly the DifferentialEquations.jl~\cite{Rackauckas2017} package and the ADTypes.jl sparse differentiation infrastructure.


\section*{Code Availability}

The complete ADORE solver source code, including all input configuration files, post-processing scripts, and the validation datasets used in this paper, is available at \url{https://github.com/YUXinluhk/BallBearing-Dynamics} under the MIT license. A permanent archival snapshot with DOI will be deposited on Zenodo upon acceptance.


%% ======================================================================
\begin{thebibliography}{99}

\bibitem{Jones1960}
A.B.~Jones, \textit{A general theory for elastically constrained ball and radial roller bearings under arbitrary load and speed conditions}, ASME J. Basic Eng. 82 (1960) 309--320.

\bibitem{Harris2006}
T.A.~Harris, M.N.~Kotzalas, \textit{Rolling Bearing Analysis}, 5th ed., CRC Press, 2006.

\bibitem{Gupta1984}
P.K.~Gupta, \textit{Advanced Dynamics of Rolling Elements}, Springer-Verlag, 1984.

\bibitem{Walters1971}
C.T.~Walters, \textit{The dynamics of ball bearings}, ASME J. Lubr. Technol. 93 (1971) 1--10.

\bibitem{Jang2004}
G.~Jang, S.W.~Jeong, \textit{Nonlinear excitation model of ball bearing waviness in a rigid rotor supported by two or more ball bearings considering five degrees of freedom}, ASME J. Tribol. 126 (2004) 82--90.

\bibitem{Bizarre2018}
L.~Bizarre, F.~Nonato, K.L.~Cavalca, \textit{Formulation of five degrees of freedom ball bearing model accounting for the nonlinear stiffness and damping of elastohydrodynamic point contacts}, Mech. Mach. Theory 124 (2018) 179--196.

\bibitem{Stacke1999}
L.E.~Stacke, D.~Fritzson, P.~Nordling, \textit{BEAST---A rolling bearing simulation tool}, Proc. Inst. Mech. Eng. Part K: J. Multi-body Dyn. 213 (1999) 63--71.

\bibitem{Palmgren1959}
A.~Palmgren, \textit{Ball and Roller Bearing Engineering}, 3rd ed., Burbank, 1959.

\bibitem{Rackauckas2017}
C.~Rackauckas, Q.~Nie, \textit{DifferentialEquations.jl---A performant and feature-rich ecosystem for solving differential equations in Julia}, J. Open Res. Softw. 5 (2017) 15.

\bibitem{Revels2016}
J.~Revels, M.~Lubin, T.~Papamarkou, \textit{Forward-mode automatic differentiation in Julia}, arXiv:1607.07892 (2016).

\bibitem{Parker1984}
R.J.~Parker, \textit{Comparison of predicted and experimental thermal performance of angular-contact ball bearings}, NASA Technical Paper 2275, 1984.

\bibitem{SKF2018}
SKF Group, \textit{Rolling Bearings}, SKF Publication PUB BU/P1 10000/2 EN, 2018.

\bibitem{Takabi2013}
J.~Takabi, M.M.~Khonsari, \textit{Experimental testing and thermal analysis of ball bearings}, Tribol. Int. 60 (2013) 93--103.

\bibitem{Bair2006}
S.~Bair, \textit{High Pressure Rheology for Quantitative Elastohydrodynamics}, Elsevier, 2007.

\bibitem{Akamatsu2009}
Y.~Akamatsu, M.~Maeda, H.~Murase, \textit{Oil supply from grooves to ball-outer-ring contacts in an angular contact ball bearing}, Tribol. Trans. 52 (2009) 636--643.

\end{thebibliography}


%% ======================================================================
\appendix
\section{Complete Parameter Tables}
\label{sec:appendix_params}

To facilitate independent reproduction, \cref{tab:params_7210b,tab:params_nasa} list all model parameters for the two validation bearings.

\begin{table}[htbp]
    \centering
    \caption{Complete model parameters for 7210B validation (Cases 1--5).}
    \label{tab:params_7210b}
    \small
    \begin{tabular}{@{}llcc@{}}
        \toprule
        Category & Parameter & Symbol & Value \\
        \midrule
        \multirow{4}{*}{Material}
        & Young's modulus & $E$ & 208~GPa \\
        & Poisson's ratio & $\nu$ & 0.3 \\
        & Ball density & $\rho_b$ & 7800~kg/m$^3$ \\
        & CTE & $\alpha_{\text{CTE}}$ & $12 \times 10^{-6}$~K$^{-1}$ \\
        \midrule
        \multirow{4}{*}{Traction}
        & Low-SRR coefficient & $\kappa_0$ & 0.08 \\
        & High-SRR coefficient & $\kappa_\infty$ & 0.04 \\
        & Transition rate & $\kappa_m$ & 0.10 \\
        & Spin friction & $\mu_{\text{spin}}$ & 0.06 \\
        \midrule
        \multirow{4}{*}{Churning}
        & Oil density & $\rho_{\text{oil}}$ & 860~kg/m$^3$ \\
        & Oil viscosity & $\mu_{\text{oil}}$ & 0.005~Pa$\cdot$s \\
        & Fill fraction & $f_{\text{fill}}$ & 0.02 \\
        & Cage web thickness & $w_{\text{cage}}$ & 0.5~mm \\
        \midrule
        \multirow{6}{*}{Thermal}
        & Ambient temperature & $T_{\text{amb}}$ & 300~K \\
        & Initial temperature & $T_{\text{init}}$ & 350~K \\
        & Specific heat (steel) & $c_{\text{steel}}$ & 460~J/(kg$\cdot$K) \\
        & $G_{\text{IR-ball}}$ & & 50~W/K \\
        & $G_{\text{OR-ball}}$ & & 50~W/K \\
        & $G_{\text{ball-oil}}$ & & 20~W/K \\
        \midrule
        \multirow{2}{*}{Damping}
        & Modal damping ratio & $\zeta$ & 0.10 \\
        & Structural damping & $c_{\text{struct}}$ & 0 \\
        \midrule
        \multirow{3}{*}{Integrator}
        & Relative tolerance & & $10^{-4}$ \\
        & Absolute tolerance & & $10^{-7}$ \\
        & Max step size & $h_{\max}$ & $10^{-4}$~s \\
        \bottomrule
    \end{tabular}
\end{table}

\begin{table}[htbp]
    \centering
    \caption{Complete model parameters for NASA 35-mm validation. All parameters not listed are identical to \cref{tab:params_7210b}.}
    \label{tab:params_nasa}
    \small
    \begin{tabular}{@{}llcc@{}}
        \toprule
        Category & Parameter & Symbol & Value \\
        \midrule
        \multirow{4}{*}{Geometry}
        & Ball diameter & $d$ & 7.14~mm \\
        & Pitch diameter & $d_m$ & 46.5~mm \\
        & Number of balls & $Z$ & 16 \\
        & Contact angle & $\alpha_0$ & 24$^\circ$ \\
        \midrule
        \multirow{3}{*}{Traction}
        & $\kappa_0$ (at 394~K) & & 0.010 \\
        & $\kappa_\infty$ (at 394~K) & & 0.005 \\
        & $\kappa_m$ & & 0.10 \\
        \midrule
        \multirow{2}{*}{Churning}
        & Fill fraction & $f_{\text{fill}}$ & 0.002 \\
        & Cage web & $w_{\text{cage}}$ & 0.3~mm \\
        \midrule
        \multirow{5}{*}{Thermal}
        & Oil inlet temp. & $T_{\text{oil,in}}$ & 394~K \\
        & Oil flow rate & $\dot{V}$ & 760~cm$^3$/min \\
        & $G_{\text{IR-ball}}$ & & 50~W/K \\
        & $G_{\text{ball-oil}}$ & & 50~W/K \\
        & $G_{\text{OR-amb}}$ & & 5~W/K \\
        \bottomrule
    \end{tabular}
\end{table}

\end{document}
